\begin{proposition}\label{prop:nonic-no115}
An elliptic curve over a number field of degree $9$ has no point of order $115$.
\end{proposition}
\begin{proof}
Put $J=J_1(115)$ and $A=(\langle16\rangle-1)J$.
The projective order of $16$ modulo $115$ is $11$.
An exact computation identifies $A$ with the isogeny factors from the
three character orbits of orders $44$, $11$ and $22$.
Each has relative cuspidal dimension $10$, and
\[
 \dim A=10\bigl(\varphi(44)+\varphi(11)+\varphi(22)\bigr)=400.
\]
All three components have full winding submodule, by the complete
level-$115$ winding record. Thus $A(\mathbb Q)$ is finite by Kato's theorem.

We show that its order is odd. The variety $A$ has good reduction at $3$.
For each of the three characters $\chi$, form the exact determinant
\[
 h_\chi=\det\bigl((1+3\chi(3))I-T_3\mid S_\chi\bigr),
\]
where $S_\chi$ is the complete sign-one cuspidal modular-symbol space,
including all oldforms. Its matrices are defined over
$L=\mathbb Q(\zeta_{44})$. The following table records the norm and
its relation to the finite reduction order of the corresponding
character component $A_\chi$:
\[
\begin{array}{r|r|r|r|c}
 \operatorname{ord}\chi&\dim_L S_\chi&[L:\mathbb Q]&
 [L:\mathbb Q(\chi)]&N_{L/\mathbb Q}(h_\chi)\bmod2\\\hline
 44&10&20&1&1\\
 11&10&20&2&1\\
 22&10&20&2&1
\end{array}
\]
The characteristic polynomial of Frobenius gives
\[
 N_{L/\mathbb Q}(h_\chi)
 =\#A_\chi(\mathbb F_3)^{[L:\mathbb Q(\chi)]}.
\]
Every norm is odd, so every $\#A_\chi(\mathbb F_3)$ is odd.
The finite reduction order is invariant under isogeny; hence
$\#A(\mathbb F_3)$ is odd. Injectivity of reduction on $2$-primary
torsion gives $A(\mathbb Q)[2^\infty]=0$.
In particular, $(\langle16\rangle-1)J(\mathbb Q)$ is finite of odd order.

We can therefore apply \cite[Proposition~4.4]{DN} at $q=2$.
Its required hypothesis on the $2$-primary part is satisfied.
Here $k_2=2(2+1)=6$, and
\[
 9k_2=54<\frac{325\cdot6336}{32768}
       =\frac{32175}{512}
       \leq\operatorname{gon}_{\mathbb C}X_1(115).
\]
Thus the trace divisor of a hypothetical degree-nine torsion point
is stable under $\langle16\rangle$. The least degree of a CM point
on $X_1(115)$ is $88$, so its support has no CM points.
The free action of this order-$11$ diamond would force $11\mid9$,
a contradiction.

The three exact determinants, their defining field polynomial and
norms are preserved in \texttt{C115\_norm\_order*.json}.
The computation \texttt{C115\_image\_even\_orders\_certified.m}
is replayed independently by \texttt{replay\_C115\_norms.py}, which
computes each norm as an exact polynomial resultant and checks its parity.
The larger integral-homology diagnostic is not used in this proof.
\end{proof}
